\section{Introdução}\label{sec:intro}

\subsection{Colapso crítico}

Em 1993, Choptuik estudou numericamente o colapso gravitacional de um campo escalar sem massa esfericamente simétrico
e encontrou, no limiar entre a dispersão e a formação de buraco negro, uma solução universal com três características
marcantes~\cite{Choptuik1993}: ela é aproximada, em todas as famílias a um parâmetro de dados iniciais que foram evoluídas, quando o
parâmetro é ajustado ao limiar; ela é discretamente autossimilar, reproduzindo-se em escalas cada vez menores com um período de eco
$\Delta\approx3.44$; e, perto do limiar, a massa do buraco negro escala como $(p-p_*)^\gamma$ com um
expoente universal $\gamma\approx0.37$. Esses \emph{fenômenos críticos} foram desde então encontrados em muitos modelos de matéria, e são
revisados em~\cite{GundlachMartinGarcia2007}. Para o sistema de Einstein--campo escalar esfericamente simétrico, Christodoulou provou que
dados suficientemente pequenos se dispersam e dados suficientemente grandes formam superfícies aprisionadas~\cite{Christodoulou1986,
Christodoulou1991,Christodoulou1993}.

A explicação padrão é dinâmica~\cite{KoikeHaraAdachi1995,Gundlach1997}. A solução crítica é um ponto
fixo (no caso discretamente autossimilar, uma órbita periódica) da evolução no tempo logarítmico, com uma variedade
estável de codimensão um, que é o próprio limiar, e um \emph{único} modo instável, que cresce como
$e^{\lambda T}$. Uma solução quase crítica fica perto da crítica até que o modo instável cresça até um
tamanho fixo, e a análise dimensional dá então $\gamma=1/\lambda$, com uma modulação periódica no caso discretamente
autossimilar. Gundlach calculou numericamente as perturbações da solução de Choptuik e encontrou exatamente um modo que cresce,
com $\gamma=0.374\pm0.001$~\cite{Gundlach1997}, e Mart\'in-Garc\'ia e Gundlach encontraram numericamente que todas as
perturbações não esféricas decaem~\cite{MartinGarciaGundlach1999}. Um ingrediente necessário desse quadro é que a
solução crítica tenha exatamente um modo instável; é esse enunciado espectral que provamos, para perturbações
esfericamente simétricas.

\subsection{Existência, e a questão espectral em aberto}

A existência do espaço-tempo de Choptuik foi provada por Reiterer e Trubowitz~\cite{RT2019}, com um argumento
assistido por computador na tradição da prova de Lanford das conjecturas de Feigenbaum~\cite{Lanford1982}. Eles constroem uma
solução analítica real e discretamente autossimilar numa vizinhança do cone de luz passado da singularidade, como
ponto fixo de uma contração perto de dados numéricos de alta precisão, e determinam a sua constante de eco e um
segundo invariante com $80$ dígitos. Sobre as perturbações, eles escrevem: ``We do not study perturbations about the
Choptuik spacetime, that are relevant for many of the (conjectured) properties of the solution. Our paper can be
the starting point for a rigorous investigation of this kind. This would be interesting, because current
numerical results seem to be inconclusive''~\cite[\S1]{RT2019} (``Não estudamos perturbações do espaço-tempo de
Choptuik, que são relevantes para muitas das propriedades (conjecturadas) da solução. O nosso artigo pode ser o ponto
de partida para uma investigação rigorosa desse tipo. Isso seria interessante, porque os resultados numéricos atuais
parecem inconclusivos'').

Este artigo faz uma investigação desse tipo para perturbações esfericamente simétricas.

\subsection{O resultado}

Trabalhamos com a solução de~\cite{RT2019}, nas suas coordenadas $(\tau,\xi)$ e variáveis $\omega$, nas quais a
autossimilaridade é a translação $\tau\mapsto\tau+2\pi$. As perturbações lineares do tipo de Floquet,
$e^{s\tau}$ vezes um perfil periódico, são regidas por uma família de operadores $L(s)$, a linearização do
sistema de Reiterer--Trubowitz, junto com uma aplicação de vínculos (\emph{constraints}) linearizada $C(s)$. O expoente $s$ corresponde à
taxa de crescimento $\lambda=2\pi s/K$ no tempo logarítmico usual, onde $e^{-K}$ é o fator pelo qual a
autossimilaridade reescala os comprimentos.

\begin{introthm}[versão informal do Teorema Principal, \cref{sec:main}]
Sob perturbações lineares esfericamente simétricas, suaves no tempo e analíticas na variável radial até o cone de luz
e através dele, e módulo gauge (calibre), o espaço-tempo de Choptuik tem exatamente um modo instável. O modo é real e algebricamente simples, e a sua taxa de crescimento satisfaz
\[
  \lambda\in[2.674040,\ 2.674115],\qquad 1/\lambda\in[0.373955,\ 0.373966].
\]
Mais precisamente, por classe de Floquet, o operador $L$ tem exatamente quatro raízes no semiplano $\Re s>0$, cada uma simples
e com um representante real: um modo de gauge, que é a translação do instante da singularidade; duas raízes cujos autovetores
violam os vínculos; e o modo físico instável. No eixo imaginário, $L$ tem apenas a raiz $s=0$, a
translação da fase. Se o cone de luz da singularidade é mantido fixo, o modo de gauge deixa de ser removido pelo
grupo menor de transformações de gauge, e a contagem passa a ser dois; a classe adicional corresponde a transladar
o instante da singularidade, e não a uma segunda instabilidade.
\end{introthm}

O valor de $1/\lambda$ concorda com o expoente crítico medido e calculado; esta é uma conferência externa, já que
a relação $\gamma=1/\lambda$ diz respeito à evolução não linear. O teorema é espectral: ele localiza as raízes de
$L$ e identifica as físicas, mas não diz nada sobre o semigrupo da evolução linearizada nem sobre a dinâmica não
linear. Ele trata apenas de perturbações esfericamente simétricas, na classe de regularidade enunciada acima. Essas
limitações são discutidas na \cref{sec:discussion}.

\subsection{Estratégia}

A prova tem três partes.

\emph{Do problema geométrico ao operador} (\cref{sec:physical}). Um modo físico é uma solução das equações
de Einstein--campo escalar linearizadas do tipo de Floquet, e só a sua classe módulo transformações de gauge tem
significado. Mostramos que todo modo desses com $\Re s>0$ pode ser posto na forma de~\cite{RT2019} por uma transformação
de gauge, e então corresponde a um elemento de $\ker L(s)\cap\ker C(s)$; que, reciprocamente, todo elemento desses
vem de um modo físico; e que os únicos modos de gauge com $\Re s>0$ são os múltiplos de um modo explícito em
$s=\mu$, a translação infinitesimal do ponto singular. A completude do gauge é global no domínio
de~\cite{RT2019} e usa que o cone de luz está contido nele. A multiplicidade física de uma raiz é então uma questão
de álgebra linear no seu espaço de raízes.

\emph{Contagem} (\cref{sec:counting}). As raízes de $L$ com parte real grande são excluídas por uma cota da
imagem numérica (\emph{numerical range}), calculada num espaço de funções no qual o fundo (\emph{background}, a solução
em torno da qual se lineariza) age ponto a ponto. Na região limitada
restante, as raízes são contadas com uma versão para operadores do teorema de Rouch\'e, devida a Gohberg e
Sigal~\cite{GohbergSigal1971}: após uma deflação (\emph{deflation}) de posto dois, que afasta as duas raízes de gauge, $L$ é comparado ao longo da
fronteira com uma referência bloco-diagonal cujos zeros são os autovalores de uma matriz $14016\times14016$, e estes
são contados a partir de uma decomposição de Schur com resíduo certificado.

\emph{Identificação} (\cref{sec:roots}). Cada raiz é localizada rigorosamente numa bola certificada (\emph{enclosure}), e mostra-se que é
algebricamente simples, por um argumento de Newton--Kantorovich aplicado a um sistema aumentado (\emph{bordered system}):
a equação de autovalor acrescida de uma condição de normalização, com o autovalor como incógnita. Mostra-se que os autovetores em duas delas violam os vínculos, por cotas inferiores
explícitas. Para a raiz restante, os vínculos não podem ser conferidos num autovetor aproximado; em vez disso, uma
identidade $KC=ML$ propaga os vínculos, e um certificado à parte mostra que o operador de propagação dos vínculos
$K$ é injetivo ali.

Todos os cálculos são feitos em aritmética racional exata, em aritmética de bolas (\emph{ball arithmetic}), ou como
cotas a posteriori sobre cálculos em ponto flutuante. Na aritmética de bolas, cada número real é representado por um
centro $m$ e um raio $r$, isto é, pelo intervalo $[m-r,m+r]$, e cada operação devolve uma bola que contém o resultado
exato. Eles levam alguns dias em computadores comuns.

\subsection{O que há de novo}

As ferramentas são clássicas; o que é novo é a sua combinação para este operador, e os resultados. O artigo transforma em
teorema um quadro espectral que só era conhecido numericamente: ele dá uma contagem rigorosa dos modos instáveis do
espaço-tempo de Choptuik; uma localização rigorosa da
taxa de crescimento física; a classificação de todas as raízes com parte real positiva; e a passagem do problema
geométrico ao operador, incluindo a completude do gauge. Os certificados, um programa de verificação que
os refaz, os dados e um guia de reprodução estão disponíveis publicamente (\cref{sec:computation-data}).

\subsection{Organização}

A \cref{sec:setting} recorda a formulação de~\cite{RT2019} e define $L$, $C$ e $K$. A \cref{sec:main} enuncia o
Teorema Principal. A \cref{sec:physical} relaciona os modos físicos ao operador, a \cref{sec:counting} conta as raízes e
a \cref{sec:roots} as identifica e completa a prova. A \cref{sec:computation} descreve os cálculos, a sua
verificação e o uso de assistentes de IA, e a \cref{sec:discussion} discute as limitações e o problema não
esférico. Os apêndices contêm as provas dos lemas estruturais, fórmulas explícitas e a lista dos certificados.

\subsection{Terminologia}\label{sec:intro-terms}

Mantivemos em português a maior parte dos termos técnicos. Na primeira ocorrência, damos entre parênteses o
equivalente em inglês, que é o que se encontra na bibliografia. A \cref{tab:termos} reúne os termos menos usuais e o
sentido em que são usados aqui.

\begin{table}[tbp]
\centering
\small
\begin{tabular}{@{}>{\raggedright\arraybackslash}p{0.22\textwidth}>{\raggedright\arraybackslash}p{0.2\textwidth}
  >{\raggedright\arraybackslash}p{0.5\textwidth}@{}}
\toprule
termo & em inglês & sentido neste artigo \\
\midrule
aritmética de bolas & \emph{ball arithmetic} & cada número real é representado por um centro $m$ e um raio $r$, isto é,
  pelo intervalo $[m-r,m+r]$; cada operação devolve uma bola que contém o resultado exato \\
célula (de uma cobertura) & \emph{tile} & cada uma das regiões de uma cobertura em que se verifica uma desigualdade; as
  células podem se sobrepor \\
deflação & \emph{deflation} & perturbação de posto finito que afasta raízes já conhecidas (aqui, as de gauge) \\
espaço do toro & \emph{torus space} & espaço $\ell^2$ ponderado de coeficientes, igual ao $L^2$ de um toro complexo, no
  qual o fundo age ponto a ponto (\cref{sec:setting-realizations}) \\
faixa A, faixa B & \emph{A band, B band} & regiões do plano do expoente $s$ que contêm as raízes dos setores A e B
  (\cref{sec:setting-sectors}) \\
fundo & \emph{background} & a solução de Reiterer e Trubowitz, em torno da qual se lineariza \\
gauge (calibre) & \emph{gauge} & liberdade de coordenadas; uma transformação de gauge soma à perturbação a derivada de
  Lie do fundo ao longo de um campo vetorial (\cref{def:gauge}) \\
imagem numérica & \emph{numerical range} & o conjunto $W(T)=\{\langle Tx,x\rangle:\ x\in\Dom(T),\ \|x\|=1\}$, onde $\Dom(T)$
  é o domínio de $T$ \\
janela, parte distante & \emph{window, far part} & os níveis intermediário e externo da referência usada na contagem
  (\cref{sec:counting-reference}) \\
localização rigorosa; bola ou intervalo certificado & \emph{enclosure} & um conjunto que comprovadamente contém o valor
  exato \\
recuperação do raio & \emph{radius recovery} & o argumento que leva um autovetor de um espaço de raio menor para o de
  raio maior (\cref{lem:recovery}) \\
setor A, setor B & \emph{sector A, sector B} & os perfis com a simetria do fundo e os com a simetria oposta
  (\cref{sec:setting-sectors}) \\
sistema aumentado & \emph{bordered system} & a equação de autovalor acrescida de uma condição de normalização, com o
  autovalor como incógnita (\cref{sec:roots-nk}) \\
testemunho & \emph{witness} & uma cota inferior que prova que um autovetor viola os vínculos (\cref{prop:witness}) \\
vínculos & \emph{constraints} & as equações do sistema que ficam fora da parte de evolução: $F^\sharp=0$ no sistema de
  Reiterer e Trubowitz e $C(s)h=0$ na linearização (\cref{sec:setting-system}) \\
\bottomrule
\end{tabular}
\caption{Termos técnicos usados neste artigo.}
\label{tab:termos}
\end{table}
